% ChargeShield-FL --- DSN 2027 --- 1. Introduction
% Scritta il 2026-09-23 (struttura a sette sezioni, guida sez. 10).
% Recepisce le note dell'utente nella versione Overleaf di
% sections/introduction.tex: definire brevemente la MIA e raccontarne i
% contributi in forma narrativa; spiegare subito la tecnica canary; nota
% sull'AUC; riformulare le frasi sui tre studi applicativi; dire che
% scambiare "l'attacco non rileva" per "non c'e' esposizione" e' pericoloso
% in esercizio; eliminare il paragrafo Purdue (troppo lungo).
% Aggiunge gli incidenti reali richiesti dall'utente (dettaglio in sez. 3.1).

\section{Introduction}
\label{sec:introduction}

A charging session is a small record: when a vehicle was plugged in, at which
charging point, for how long, and how much energy it drew. Linked to an
account or to a vehicle identifier, a sequence of such records describes a
person's routine. The systems that hold these records have repeatedly failed
to keep them. In June 2023 a Shell Recharge Solutions database, reachable
from the Internet without a password, exposed logging data containing fleet
customers' names and contacts, vehicle identification numbers (VINs) and the
locations of charging points, residential ones
included~\cite{whittaker2023shell}. In November 2024 about 116{,}000 records
from several charge-point operators that shared a third-party charging
application appeared on a hacking forum; according to the published analysis
they contained names, addresses, vehicle models and VINs, precise
geolocation and the OCPP logs exchanged between operators and
chargers~\cite{robb2024cpo}. In March 2026 a ransomware attack on the cloud
platform of the charger manufacturer ELECQ encrypted and copied customer
databases with names, e-mail and home addresses and phone numbers, and was
notified to data-protection authorities in the United Kingdom and
Germany~\cite{page2026elecq}. Section~\ref{sec:threat-incidents} summarises
these incidents, and what they do and do not say about the threat studied
here.

Federated learning (FL)~\cite{mcmahan2017communication} is often proposed
for such settings because the records stay with the operator that collected
them: each site trains a model locally and shares only an update, which an
aggregator combines into a global model. Not sharing the records, however, is
not the same as not revealing them. A \emph{membership inference attack}
(MIA) asks, for a given record and a trained model or update, whether that
record was part of the training data. The question looks modest, but for a
charging session a positive answer places a specific vehicle at a specific
site at a specific time, and it is the elementary test on which most
empirical evaluations of training-data leakage are built.

The attack has a short and instructive history. Shokri
et~al.~\cite{shokri2017membership} showed that membership leaks through a
model's outputs by training \emph{shadow models} that imitate the target on
data whose membership the attacker knows, and learning from them how members
and non-members differ. Yeom et~al.~\cite{yeom2018privacy} then observed that
much of that signal is already in the loss: a model fits its training records
better than unseen ones, so thresholding the loss is itself an attack, and its
success grows with overfitting. Salem et~al.~\cite{salem2019mlleaks} relaxed
the attacker's knowledge requirements, showing that a single shadow model
trained on different data often suffices. Carlini
et~al.~\cite{carlini2022membership} recast the problem as a hypothesis test
on each record --- the likelihood-ratio attack (LiRA), which compares a
record's loss with the losses it obtains in shadow models trained with and
without it --- and argued that an attack should be judged by how many members
it identifies at a very low false-positive rate, not by its average accuracy.
In FL the same question can be put to the global model, which every
participant receives, or to an individual update, which the aggregator
receives~\cite{nasr2019comprehensive}.

Differential privacy (DP)~\cite{dwork2006calibrating,dwork2014algorithmic}
is the standard answer to this risk: noise calibrated to a bounded
contribution makes the output nearly insensitive to the presence of any one
protected unit, whatever the attacker does. The guarantee, however, is stated
with respect to a unit, a mechanism and its outputs. In FL the unit can be a
whole client~\cite{mcmahan2018learning} or a single record~\cite{abadi2016deep},
and the two protect different things. The parameter $\varepsilon$ is not a
fraction of identifiable records, and a weak attack shows neither that the
protection is effective nor that it is useless. Empirical evaluation answers
a complementary question: what \emph{specified} attacks can infer under the
observed conditions~\cite{jayaraman2019evaluating,jagielski2020auditing}.

Reading such evaluations requires care. Attack results are usually summarised
by the area under the ROC curve (AUC)\footnote{The AUC is the probability that
the attack scores a randomly chosen member above a randomly chosen
non-member: $0.5$ is random guessing and $1$ perfect separation. Because it
averages over all decision thresholds, an AUC close to $0.5$ is compatible
with a few members being identified with high
confidence~\cite{carlini2022membership}.}, and DP is often credited with
bringing it close to $0.5$. In several applied studies, however, the attack was
already at that level without protection, or no unprotected reference is
shown. A study on cross-border telemedicine reports an AUC of about $0.50$ for
centralised, federated and differentially private federated training
alike~\cite{meyo2025telemedicine}; a study on maritime vessel data presents an
AUC of about $0.51$ under DP as evidence of strong
protection~\cite{khan2026maritime}; a study on precision medicine reports an
unprotected centralised baseline of $0.5486$~\cite{pasupuleti2025biointelligence}.
\Susy{Meyo et al.\ e Khan et al.\ verificati sulla pagina dell'editore il
2026-09-23 (Khan: non ho visto un riferimento senza DP, da controllare nel
testo completo). Pasupuleti: pagina non raggiungibile, valore 0.5486 non
verificato sul testo originale (guida sez. 3 e 11).}
When an attack cannot separate members from non-members before any protection
is applied, an AUC near $0.5$ afterwards cannot show that DP suppressed a
signal. Two hypotheses must then be kept apart: that the system exposes no
membership signal detectable by the attack, and that the attack cannot detect
a signal that is present. The difference is not only methodological. An
operator who takes the second situation for the first concludes that no
exposure exists and may deploy on that basis, while an adversary with a
better-suited attack could still succeed. A measurement that returns chance
is a statement about the instrument and its precision, not a certificate that
the system is safe.

The instrument itself can be tested with a \emph{positive control} based on
canaries~\cite{maddock2023canife,steinke2023privacy}. The experimenter inserts
known records into the training data --- here, real sessions duplicated many
times so that the model is pushed to memorise them --- and keeps an equally
chosen set of sessions out. Membership is then present by construction and
the attack should find it; if it does not, its null results elsewhere say
little. Two refinements make the control informative. Exchanging the roles of
the two sets in a second training run separates membership from intrinsic
difficulty, because a membership signal follows whichever set is trained on.
Scoring the same records on an untrained model measures how far apart the two
sets were before any training. A passed control shows that the instrument is
sensitive in the regime in which it was run; it does not certify null results
obtained in another regime.

This paper reports an empirical study of this question in a concrete system:
federated autoencoders trained on real charging sessions from the three sites
of ACN-Data~\cite{lee2019acndata}, one client per site, extending the problem
statement of our prior work~\cite{imperatrice2026toward}. Our objective is to
establish \emph{under which conditions DP reduces the ability of membership
inference attacks in this federated system, and at what cost to the model}. We
keep the research questions of the original project:

\begin{itemize}
\item \textbf{RQ1.} How do attack performance and reconstruction error vary
with the DP configuration and the number of rounds, with data, partition and
federated algorithm fixed? The protected unit, client or record, is a factor
of RQ1.
\item \textbf{RQ3.} With data, partition, training resources and a comparable
DP guarantee fixed, does FedProx with $\mu = 0.01$~\cite{li2020fedprox} change
the attack--reconstruction trade-off relative to FedAvg ($\mu = 0$)?
\item \textbf{RQ2.} With the same population and learning configuration, how
do attacks and the cost of DP change between an approximately IID partition
and the per-site partition?
\end{itemize}

\noindent We address them in the order RQ1, RQ3, RQ2, because the algorithm
comparison reuses the central protocol most directly. A fourth question of the
original project, detecting a passive attacker from anomalies in the
gradients, is excluded for the reason given in Section~\ref{sec:threat-model}.

We evaluate two DP mechanisms --- client-level clipping and noise on the
update, and record-level DP-SGD --- and, for the client-level one, three
observation points along the privatisation pipeline, ordered by the
information they leave to the adversary. We measure attacks on the global
model (loss threshold and a shadow-calibrated threshold) and on the per-client
update (LiRA), a per-record analysis of which sessions are recognised
consistently across seeds, and the reconstruction error on a fixed natural
hold-out. Every cell uses five seeds, and metrics, surfaces and thresholds were
fixed before the central sweep was read (Section~\ref{sec:method-spec}).

The evidence verified so far establishes three facts. Without DP, no attack
separates members from non-members on average, yet roughly 120 sessions above
the chance level are recognised consistently across seeds. Client-level DP at
$\varepsilon \le 1$ per round removes that excess but raises the
reconstruction error more than 100-fold, so in those cells protection and loss
of utility cannot be separated. And the measurement detects induced
memorisation on one site, symmetrically under role exchange. The experiments
that can place an operating point between these regimes --- the sweep over
$\varepsilon \in [2, 16]$ and the record-level comparison --- and the RQ3 and
RQ2 comparisons are \pending{running or pending: see Table~\ref{tab:matrix}}.

Our contributions are the following.

\begin{itemize}
\item \textbf{A controlled characterisation of the DP--membership--reconstruction
trade-off} on real multi-site EV-charging data, with two protected units and
three observation points, reported together with the cost on the model and
with per-seed variability (RQ1, Section~\ref{sec:results-rq1}).
\pending{result of the $\varepsilon$ sweep and of record-level DP}

\item \textbf{Paired comparisons} of FedAvg against FedProx (RQ3) and of an
IID against the per-site partition (RQ2), with the same records, seeds and
attacks. \pending{to be reduced explicitly if E-E or E-D are not completed
(guide, Sec.~9)}

\item \textbf{Measurement validity}: a canary control with role exchange and
untrained-model baselines, a per-record analysis with a permutation chance
level, and evidence that the shadow calibration of LiRA degenerates in this
regime, which makes the uncalibrated loss the metric we rely on.

\item \textbf{A reproducible evaluation}: configurations, seeds and analysis
scripts for a single-process simulation and for a five-container NVIDIA FLARE
deployment that runs the same attack code.
\end{itemize}

Section~\ref{sec:background} introduces DP, membership inference and the
closest related work. Section~\ref{sec:system} describes the data, the
federation, the DP mechanisms and the threat model, starting from the
incidents above. Section~\ref{sec:method} specifies the experiments and the
rules fixed before reading them, Section~\ref{sec:results} reports the
results, Section~\ref{sec:discussion} states what they allow us to conclude and
what they do not, and Section~\ref{sec:conclusion} concludes.
